% Reviewed translation: PDF pages 320--325 / printed pages 306--311.
\MathBookSourcePage{320}
\subsection{一类非线性问题的应用}
\label{subsec:navier-branch-nonlinear-application}

现在把前述理论方法用于求解如下问题类:
\begin{equation}\label{eq:navier-branch-operator-problem}
F(\lambda,u)\equiv u+TG(\lambda,u)=0,
\end{equation}
其中 $T\in\mathcal L(Y;X)$, $G$ 是从 $\Lambda\times X$ 到 $Y$ 的 $C^2$ 映射,
$X,Y$ 是两个 Banach 空间, $\Lambda$ 是 $\mathbb R$ 的紧区间. 因而此处
$\mathscr X=X$. 我们在 Subsection~\ref{subsec:navier-nonsingular-abstract}
中已经看到, Navier--Stokes 方程的 Dirichlet 问题可以写成
\eqref{eq:navier-branch-operator-problem} 的形式.

为逼近 Problem~\eqref{eq:navier-branch-operator-problem}, 引入一个用来逼近
$T$ 的算子 $T_h\in\mathcal L(Y;X)$, 并置
\begin{equation}\label{eq:navier-branch-operator-approximation}
F_h(\lambda,u)=u+T_hG(\lambda,u).
\end{equation}
逼近问题为:
\begin{equation}\label{eq:navier-branch-operator-discrete-problem}
\text{find }u_h\in X\text{ such that }F_h(\lambda,u_h)=0.
\end{equation}

现设 Problem~\eqref{eq:navier-branch-operator-problem} 有一个非奇异解分支
$\{(\lambda,u(\lambda));\lambda\in\Lambda\}$. 为应用
Theorem~\ref{thm:navier-branch-uniform-branch}, 还需要附加假设. 首先设存在连续嵌入
$Y$ 的另一个 Banach 空间 $Z$, 使得
\begin{equation}\label{eq:navier-branch-derivative-into-z}
D_uG(\lambda,u)\in\mathcal L(X;Z)
\qquad\forall\lambda\in\Lambda,\quad\forall u\in X.
\end{equation}
其次, 对算子 $T_h$ 的逼近性质作如下假设:
\begin{align}
\lim_{h\to0}\|(T_h-T)g\|_X&=0
&&\forall g\in Y,
\label{eq:navier-branch-strong-operator-convergence}\\
\lim_{h\to0}\|T_h-T\|_{\mathcal L(Z;X)}&=0.
\label{eq:navier-branch-z-operator-convergence}
\end{align}

\setcounter{statement}{3}
\begin{Remark}\label{rem:navier-branch-compact-operator-approximation}
当算子 $T_h$ 的像是 $X$ 的有限维子空间 $X_h$ 时(有限元逼近正是这种情形),
假设 \eqref{eq:navier-branch-z-operator-convergence} 蕴含 $T$ 是从 $Z$ 到 $X$
的紧算子, 因为它是紧算子序列 $T_h$ 的一致极限. 因而在这种情形下,
$TD_uG(\lambda,u)\in\mathcal L(X;X)$ 是紧的, 而
\[
D_uF(\lambda,u)=I+TD_uG(\lambda,u)
\]
是恒等算子的紧扰动.

还应注意, 当 $Z$ 到 $Y$ 的嵌入是紧的时,
\eqref{eq:navier-branch-z-operator-convergence} 是
\eqref{eq:navier-branch-strong-operator-convergence} 的推论.
\end{Remark}

\MathBookSourcePage{321}
\setcounter{statement}{2}
\begin{Theorem}\label{thm:navier-branch-operator-approximation}
设 $G$ 是从 $\Lambda\times X$ 到 $Y$ 的 $C^2$ 映射, 且映射 $D^2G$
在 $\Lambda\times X$ 的每个有界子集上有界. 再设
\eqref{eq:navier-branch-derivative-into-z}--\eqref{eq:navier-branch-z-operator-convergence}
成立, 且 $\{(\lambda,u(\lambda));\lambda\in\Lambda\}$ 是
\eqref{eq:navier-branch-operator-problem} 的一个非奇异解分支. 则存在 $X$ 中原点的
一个邻域 $\mathcal O$, 并且对充分小的 $h\leq h_0$, 存在唯一的 $C^2$ 映射
$\lambda\in\Lambda\mapsto u_h(\lambda)\in X$, 使得
\begin{align}
\{(\lambda,u_h(\lambda));\lambda\in\Lambda\}
&\text{ is a branch of nonsingular solutions of }
\eqref{eq:navier-branch-operator-discrete-problem},
\label{eq:navier-branch-operator-discrete-branch}\\
u_h(\lambda)-u(\lambda)&\in\mathcal O
\qquad\forall\lambda\in\Lambda.
\label{eq:navier-branch-operator-neighborhood}
\end{align}
此外, 存在与 $h$ 和 $\lambda$ 无关的常数 $K>0$, 使
\begin{equation}\label{eq:navier-branch-operator-error}
\|u_h(\lambda)-u(\lambda)\|_X
\leq K\|(T_h-T)G(\lambda,u(\lambda))\|_X
\qquad\forall\lambda\in\Lambda.
\end{equation}
\end{Theorem}
\begin{Proof}
以 $\widetilde u_h(\lambda)=u(\lambda)$ 应用
Theorem~\ref{thm:navier-branch-uniform-branch}. 由假设,
$(\lambda,u)\mapsto G(\lambda,u)$ 是 $C^2$ 映射, 故
$(\lambda,u)\mapsto F(\lambda,u)$ 以及 $\lambda\mapsto u(\lambda)$ 也是 $C^2$ 的.

先验证 \eqref{eq:navier-branch-uniform-mu}. 有
\[
\mu_h(\lambda)=
\|(T-T_h)D_uG(\lambda,u(\lambda))\|_{\mathcal L(X;X)}.
\]
于是 \eqref{eq:navier-branch-derivative-into-z}、
\eqref{eq:navier-branch-z-operator-convergence} 以及映射
$\lambda\mapsto D_uG(\lambda,u(\lambda))$ 的连续性给出
\eqref{eq:navier-branch-uniform-mu}.

接着验证 \eqref{eq:navier-branch-uniform-residual}. 因为
$F(\lambda,u(\lambda))=0$, 可写成
\[
\varepsilon_h(\lambda)
=\|F_h(\lambda,u(\lambda))-F(\lambda,u(\lambda))\|_X
=\|(T_h-T)G(\lambda,u(\lambda))\|_X.
\]
映射 $\lambda\mapsto G(\lambda,u(\lambda))$ 的连续性和
\eqref{eq:navier-branch-strong-operator-convergence} 因而蕴含
\eqref{eq:navier-branch-uniform-residual}.

最后考察 \eqref{eq:navier-branch-uniform-lipschitz}. 由
\eqref{eq:navier-branch-strong-operator-convergence} 和一致有界原理可得
\[
\|T_h\|_{\mathcal L(Y;X)}\leq C.
\]
因此
\[
L_h(\lambda;\alpha)
\leq C\sup_{v\in S(u(\lambda);\alpha)}
\|D_uG(\lambda,u(\lambda))-D_uG(\lambda,v)\|_{\mathcal L(X;Y)}.
\]
由中值定理,
\[
L_h(\lambda;\alpha)\leq\alpha C L(\alpha),
\qquad
L(\alpha)=
\sup_{\substack{\lambda\in\Lambda\\v\in S(u(\lambda);\alpha)}}
\|D_{uu}^2G(\lambda,v)\|_{\mathcal L_2(X;Y)}.
\]
由于 $D^2G$ 在 $\Lambda\times X$ 的每个有界子集上有界,
立即得到 \eqref{eq:navier-branch-uniform-lipschitz}.

所以 Theorem~\ref{thm:navier-branch-uniform-branch} 的结论成立.
估计 \eqref{eq:navier-branch-operator-error} 直接由
\eqref{eq:navier-branch-uniform-error} 得到, 而 $u_h$ 的 $C^2$ 正则性由
$G$ 的正则性得到, 参见 Remark~\ref{rem:navier-branch-regularity-in-parameter}.
\end{Proof}

\MathBookSourcePage{322}
\setcounter{statement}{4}
\begin{Remark}\label{rem:navier-branch-lipschitz-derivative}
除去 $u_h$ 的 $C^2$ 正则性, Theorem~\ref{thm:navier-branch-operator-approximation}
的结论也可以把 $G$ 的 $C^2$ 正则性替换为 $D_uG$ 的 Lipschitz 连续性得到:
存在局部有界函数 $\mu\mapsto L(\mu):\mathbb R_+\to\mathbb R_+$, 使对所有
$v\in S(u(\lambda);\mu)$ 及所有 $\lambda\in\Lambda$,
\begin{equation}\label{eq:navier-branch-local-lipschitz-derivative}
\|D_uG(\lambda,u(\lambda))-D_uG(\lambda,v)\|_{\mathcal L(X;Y)}
\leq L(\mu)\|u(\lambda)-v\|_X.
\end{equation}
\end{Remark}

\setcounter{statement}{5}
\begin{Remark}\label{rem:navier-branch-higher-regularity}
另一方面, 当 $G$ 是 $C^p$ 映射($p\geq2$), 且 $D^pG$ 在
$\Lambda\times X$ 的每个有界子集上有界时,
Brezzi, Rappaz \& Raviart~\cite{BrezziRappazRaviart1980} 的论证表明
$u_h(\lambda)$ 是从 $\Lambda$ 到 $X$ 的 $C^p$ 映射, 并对每个
$0\leq m\leq p-1$ 给出估计
\begin{equation}\label{eq:navier-branch-higher-derivative-error}
\left\|\frac{d^m}{d\lambda^m}(u_h(\lambda)-u(\lambda))\right\|_X
\leq C_m\sum_{l=0}^{m}
\left\|(T_h-T)\frac{d^l}{d\lambda^l}G(\lambda,u(\lambda))\right\|_X.
\end{equation}
\end{Remark}

作为 Theorem~\ref{thm:navier-branch-operator-approximation} 的第一个应用,
把 Subsection~\ref{subsec:stokes-dirichlet-velocity-pressure} 中对 Stokes 方程引入的正则化方法或罚方法
推广到 Navier--Stokes 方程 \eqref{eq:navier-steady-system}. 考虑问题
\begin{equation}\label{eq:navier-branch-penalty-system}
\left\{
\begin{aligned}
-\nu\Delta\boldsymbol u^\varepsilon
+\sum_{j=1}^{N}u_j^\varepsilon
\frac{\partial\boldsymbol u^\varepsilon}{\partial x_j}
+\operatorname{grad}p^\varepsilon&=\boldsymbol f
&&\text{in }\Omega,\\
p^\varepsilon&=-\frac1\varepsilon\operatorname{div}\boldsymbol u^\varepsilon
&&\text{in }\Omega,\\
\boldsymbol u^\varepsilon&=\boldsymbol g
&&\text{on }\Gamma,
\end{aligned}
\right.
\end{equation}
其中给定 $\varepsilon>0$, 求
$(\boldsymbol u^\varepsilon,p^\varepsilon)\in H^1(\Omega)^N\times L_0^2(\Omega)$.
等价地, 求 $\boldsymbol u^\varepsilon\in H^1(\Omega)^N$ 使
\begin{equation}\label{eq:navier-branch-penalty-velocity}
\left\{
\begin{aligned}
-\nu\Delta\boldsymbol u^\varepsilon
-\frac1\varepsilon\operatorname{grad}(\operatorname{div}\boldsymbol u^\varepsilon)
+\sum_{j=1}^{N}u_j^\varepsilon
\frac{\partial\boldsymbol u^\varepsilon}{\partial x_j}
&=\boldsymbol f &&\text{in }\Omega,\\
\boldsymbol u^\varepsilon&=\boldsymbol g &&\text{on }\Gamma.
\end{aligned}
\right.
\end{equation}

为了研究这一正则化方法的收敛性, 考虑方程
\eqref{eq:navier-steady-system}
在 $\mathbb R_+\setminus\{0\}$ 的紧区间 $\Lambda$ 上的一个非奇异解分支
\[
\{(\lambda,u(\lambda))=(\boldsymbol u(\lambda),\lambda p(\lambda));
\lambda=1/\nu\in\Lambda\}.
\]
这意味着对每个
$(\boldsymbol u,\lambda p)=(\boldsymbol u(\lambda),\lambda p(\lambda))$,
$\lambda=1/\nu\in\Lambda$, 线性化问题 \eqref{eq:navier-branch-linearized-problem}
是适定的.

\setcounter{statement}{3}
\begin{Theorem}\label{thm:navier-branch-penalty-approximation}
设 $N\leq4$, 且
$\{(\lambda,u(\lambda))=(\boldsymbol u(\lambda),\lambda p(\lambda));
\lambda=1/\nu\in\Lambda\}$ 是方程
\eqref{eq:navier-steady-system}
的一个非奇异解分支. 则存在
$H^1(\Omega)^N\times L_0^2(\Omega)$ 中原点的一个邻域 $\mathcal O$, 并且对充分小的
$\varepsilon\leq\varepsilon_0$, 存在唯一的 $C^\infty$ 解分支
\[
\{(\lambda,u^\varepsilon(\lambda))
=(\boldsymbol u^\varepsilon(\lambda),\lambda p^\varepsilon(\lambda));
\lambda\in\Lambda\}
\]
是 \eqref{eq:navier-branch-penalty-system} 的解分支, 使
\[
u^\varepsilon(\lambda)\in u(\lambda)+\mathcal O
\qquad\forall\lambda\in\Lambda.
\]
此外有估计
\begin{equation}\label{eq:navier-branch-penalty-error}
\sup_{\lambda\in\Lambda}
\left(\|\boldsymbol u^\varepsilon(\lambda)-\boldsymbol u(\lambda)\|_{1,\Omega}
+\|p^\varepsilon(\lambda)-p(\lambda)\|_{0,\Omega}\right)
\leq C\varepsilon,
\end{equation}
其中常数 $C$ 与 $\varepsilon$ 无关.
\end{Theorem}

\MathBookSourcePage{323}
\begin{Proof}
用 \eqref{eq:navier-branch-ns-spaces} 和
\eqref{eq:navier-branch-ns-data-space} 定义空间 $X,Y$, 并分别用
\eqref{eq:navier-branch-stokes-solution-operator} 和 \eqref{eq:navier-branch-ns-nonlinear-map}
定义映射 $T,G$. 我们已经看到,
Problem~\eqref{eq:navier-steady-system}
适合 Subsection~\ref{subsec:navier-nonsingular-abstract} 的框架. 为应用
Theorem~\ref{thm:navier-branch-operator-approximation}, 取 $Z=Y$, 并把
$T^\varepsilon\in\mathcal L(Y;X)$ 定义如下: 对给定
$(\boldsymbol f_*,\boldsymbol g_*)\in Y$, 令
$(\boldsymbol u_*^\varepsilon,p_*^\varepsilon)
=T^\varepsilon(\boldsymbol f_*,\boldsymbol g_*)$ 表示正则化 Stokes 问题
\[
\left\{
\begin{aligned}
-\Delta\boldsymbol u_*^\varepsilon+\operatorname{grad}p_*^\varepsilon
&=\boldsymbol f_* &&\text{in }\Omega,\\
p_*^\varepsilon&=-\frac1\varepsilon
\operatorname{div}\boldsymbol u_*^\varepsilon &&\text{in }\Omega,\\
\boldsymbol u_*^\varepsilon&=\boldsymbol g_* &&\text{on }\Gamma
\end{aligned}
\right.
\]
的解. 显然, $(\boldsymbol u^\varepsilon,p^\varepsilon)$ 是
\eqref{eq:navier-branch-penalty-system} 的解, 当且仅当
\[
u^\varepsilon+T^{\varepsilon\nu}G(\lambda,u^\varepsilon)=0,
\qquad
u^\varepsilon=(\boldsymbol u^\varepsilon,p^\varepsilon/\nu).
\]
此外, Theorem~\ref{thm:stokes-dirichlet-penalty-error} 表明, 对所有充分小的
$\varepsilon\leq\varepsilon_0$,
\[
\|T^{\varepsilon\nu}-T\|_{\mathcal L(Y;X)}
\leq C_1\varepsilon\nu\leq C_2\varepsilon,
\]
其中由于 $\Lambda$ 的紧性, 常数 $C_2$ 与 $\lambda$ 无关.

又因为
\[
G(\lambda,v)=
\left(\lambda\left(\sum_{j=1}^{N}v_j
\frac{\partial\boldsymbol v}{\partial x_j}-\boldsymbol f\right),
-\boldsymbol g\right),
\]
$D^2G$ 与 $v$ 无关:
\[
D^2G(\lambda,v)\cdot(u,w)
=\lambda\sum_{j=1}^{N}\left(
w_j\frac{\partial u}{\partial x_j}
+u_j\frac{\partial w}{\partial x_j}\right).
\]
由 Sobolev 嵌入 Theorem~\ref{thm:sobolev-embedding} 可知, 当 $N\leq4$ 时,
$D^2G$ 在 $\Lambda\times X$ 的每个有界子集上有界.
(更一般地, $G$ 是 $C^\infty$ 的, 且对所有 $p\geq2$, $D^pG=0$.)
因此, $\{(\lambda,u(\lambda));\lambda\in\Lambda\}$ 是
\[
F(\lambda,u)\equiv u+TG(\lambda,u)=0
\]
的非奇异解分支这一事实允许应用
Theorem~\ref{thm:navier-branch-operator-approximation}. 换言之, 若
$\varepsilon_0$ 足够小, 则存在实数 $a>0$ 以及问题
\[
F^\varepsilon(\lambda,u^\varepsilon)
\equiv u^\varepsilon+T^{\varepsilon\nu}G(\lambda,u^\varepsilon)=0,
\qquad \nu=1/\lambda,
\]
的唯一非奇异解分支
$\{(\lambda,u^\varepsilon(\lambda))
=(\boldsymbol u^\varepsilon(\lambda),\lambda p^\varepsilon(\lambda));
\lambda\in\Lambda\}$, 满足
\[
\|u^\varepsilon(\lambda)-u(\lambda)\|_X\leq a.
\]
并且
\begin{align*}
\|u^\varepsilon(\lambda)-u(\lambda)\|_X
&=\|\boldsymbol u^\varepsilon(\lambda)-\boldsymbol u(\lambda)\|_{1,\Omega}
+\lambda\|p^\varepsilon(\lambda)-p(\lambda)\|_{0,\Omega}\\
&\leq C_3\|(T^{\varepsilon\nu}-T)G(\lambda,u(\lambda))\|_X\\
&\leq C_2C_3\varepsilon\|G(\lambda,u(\lambda))\|_X
\leq C_4\varepsilon.
\end{align*}
再由 Remark~\ref{rem:navier-branch-higher-regularity}, 映射
$\lambda\mapsto u^\varepsilon(\lambda)$ 是从 $\Lambda$ 到 $X$ 的 $C^\infty$ 映射.
\end{Proof}

\MathBookSourcePage{324}
令 $H$ 是满足 $X\hookrightarrow H$ 的 Banach 空间, 其中像往常一样,
符号 $\hookrightarrow$ 表示嵌入连续. 现在希望导出
$\|u_h(\lambda)-u(\lambda)\|_H$ 的更精细估计. 为此, 假设存在另一个满足
$W\hookrightarrow X$ 的 Banach 空间 $W$, 并且有如下性质:
\begin{equation}\label{eq:navier-branch-extended-derivative}
\left\{
\begin{aligned}
&\text{for all }v\in W,\ D_uG(\lambda,v)\text{ may be extended as an operator in }\mathcal L(H;Y),\\
&v\mapsto D_uG(\lambda,v)\text{ being continuous from }W\text{ into }\mathcal L(H;Y).
\end{aligned}
\right.
\end{equation}
因而对 $v\in W$, $D_uF(\lambda,v)$ 和 $D_uF_h(\lambda,v)$ 都可以延拓为
$\mathcal L(H;H)$ 的算子. 再假设
\begin{equation}\label{eq:navier-branch-h-operator-convergence}
\lim_{h\to0}\|T_h-T\|_{\mathcal L(Y;H)}=0.
\end{equation}

\setcounter{statement}{6}
\begin{Remark}\label{rem:navier-branch-compact-embedding-h}
当 $X$ 到 $H$ 的嵌入是紧的时,
\eqref{eq:navier-branch-h-operator-convergence} 同样是
\eqref{eq:navier-branch-strong-operator-convergence} 的推论.
\end{Remark}

于是可以证明如下结果.

\setcounter{statement}{4}
\begin{Theorem}\label{thm:navier-branch-sharper-norm-error}
保留 Theorem~\ref{thm:navier-branch-operator-approximation} 的假设, 并再加上
\eqref{eq:navier-branch-extended-derivative} 和
\eqref{eq:navier-branch-h-operator-convergence}. 另外假设:
\begin{align}
&\text{for each }\lambda\in\Lambda,\quad
u(\lambda)\in W\text{ and }\lambda\mapsto u(\lambda)
\in C^0(\Lambda;W),
\label{eq:navier-branch-solution-in-w}\\
&\text{for each }\lambda\in\Lambda,\quad
D_uF(\lambda,u(\lambda))\text{ is an isomorphism of }H.
\label{eq:navier-branch-derivative-isomorphism-h}
\end{align}
则对充分小的 $h\leq h_1$, 存在与 $h,\lambda$ 无关的常数 $K'>0$, 使
\begin{equation}\label{eq:navier-branch-sharper-error}
\|u_h(\lambda)-u(\lambda)\|_H
\leq K'\left\{
\|(T-T_h)G(\lambda,u(\lambda))\|_H
+\|u_h(\lambda)-u(\lambda)\|_X^2
\right\}.
\end{equation}
\end{Theorem}

\MathBookSourcePage{325}
\begin{Proof}
首先由 \eqref{eq:navier-branch-extended-derivative} 和
\eqref{eq:navier-branch-solution-in-w} 可见
\begin{align*}
&\|D_uF(\lambda,u(\lambda))-D_uF_h(\lambda,u(\lambda))\|_{\mathcal L(H;H)}\\
&\qquad=
\|(T-T_h)D_uG(\lambda,u(\lambda))\|_{\mathcal L(H;H)}\\
&\qquad\leq
\|T-T_h\|_{\mathcal L(Y;H)}
\|D_uG(\lambda,u(\lambda))\|_{\mathcal L(H;Y)}.
\end{align*}
因此 \eqref{eq:navier-branch-h-operator-convergence} 蕴含
\[
\lim_{h\to0}\sup_{\lambda\in\Lambda}
\|D_uF(\lambda,u(\lambda))-D_uF_h(\lambda,u(\lambda))\|_{\mathcal L(H;H)}=0.
\]
由 \eqref{eq:navier-branch-derivative-isomorphism-h} 和
Lemma~\ref{lem:navier-branch-derivative-perturbation}, 对所有充分小的 $h$
以及所有 $\lambda\in\Lambda$, $D_uF_h(\lambda,u(\lambda))$ 是 $H$ 上的同构, 且
\[
\|[D_uF_h(\lambda,u(\lambda))]^{-1}\|_{\mathcal L(H;H)}\leq C_1.
\]
然后像在 Theorem~\ref{thm:navier-branch-local-existence-error} 中一样可写成
\begin{align*}
u(\lambda)-u_h(\lambda)
=[D_uF_h(\lambda,u(\lambda))]^{-1}\big[&T_h\{D_uG(\lambda,u(\lambda))
\cdot(u(\lambda)-u_h(\lambda))\\
&-(G(\lambda,u(\lambda))-G(\lambda,u_h(\lambda)))\}
+F_h(\lambda,u(\lambda))\big].
\end{align*}
而
\begin{equation}\label{eq:navier-branch-taylor-remainder}
\begin{aligned}
&G(\lambda,u(\lambda))-G(\lambda,u_h(\lambda))
-D_uG(\lambda,u(\lambda))\cdot(u(\lambda)-u_h(\lambda))\\
&\qquad=-\int_0^1(1-t)D_{uu}^2G
(\lambda,(1-t)u(\lambda)+tu_h(\lambda)),dt
\cdot(u(\lambda)-u_h(\lambda))^2.
\end{aligned}
\end{equation}
因此 $D^2G$ 的有界性假设与
\eqref{eq:navier-branch-h-operator-convergence} 一起给出
\[
\|u(\lambda)-u_h(\lambda)\|_H
\leq C_1\left[
C_2\|u(\lambda)-u_h(\lambda)\|_X^2
+\|F_h(\lambda,u(\lambda))-F(\lambda,u(\lambda))\|_H
\right],
\]
这就证明了 \eqref{eq:navier-branch-sharper-error}.
\end{Proof}
